\documentclass[11pt]{article}

\usepackage[T1]{fontenc}
\usepackage{amsmath,amssymb}
\usepackage{booktabs,tabularx,array}
\usepackage[margin=1in]{geometry}
\usepackage{microtype}
\usepackage{enumitem}
\usepackage{xcolor}
\usepackage{hyperref}
\usepackage{titlesec}

\definecolor{irisblue}{HTML}{22577A}
\definecolor{irisgray}{HTML}{4B5563}
\hypersetup{
  colorlinks=true,
  linkcolor=irisblue,
  urlcolor=irisblue,
  pdftitle={A Self-Gated Adaptive Hill-q Rule for Long-Peptide Aggregation},
  pdfauthor={IRIS methodological development memo}
}

\titleformat{\section}{\Large\bfseries\color{irisblue}}{\thesection}{0.65em}{}
\titleformat{\subsection}{\large\bfseries\color{irisgray}}{\thesubsection}{0.65em}{}
\setlength{\parindent}{0pt}
\setlength{\parskip}{0.65em}
\setlist{itemsep=0.25em,topsep=0.35em}
\renewcommand{\arraystretch}{1.16}

\newcommand{\LSE}{\operatorname{LSE}}
\newcommand{\sigmoid}{\operatorname{sigmoid}}

\title{\textbf{A Self-Gated Adaptive Hill-$q$ Rule for\\Long-Peptide Aggregation}}
\author{IRIS methodological development memo}
\date{August 19, 2026}

\begin{document}
\maketitle

\begin{abstract}
IRIS scores candidate minimal-epitope--HLA pairs, whereas the PDAC and
SARS-CoV-2 assays label unresolved long peptides. A Level~2 rule must therefore
combine a variable roster of candidate scores into one long-peptide score.
Maximum aggregation represents one-candidate sufficiency but ignores
corroboration. Log-sum-exp accumulates corroboration but can reward candidate
opportunity. The proposed self-gated Hill-$q$ rule lies between these
endpoints. It starts at the strongest candidate, forms a score-weighted
corroboration offer from the roster, and admits only as much of that offer as
the resulting aggregate still needs. The Hill order $q$ controls how candidate
breadth is recognized. The parameters $(q,c,\kappa)$ may be calibrated for a
component model and metric branch, but they remain fixed across the PDAC and
SARS-CoV-2 contexts to which that scoring system is applied.
\end{abstract}

\section{Why a Level 2 aggregation rule is needed}

\subsection{The experimental unit is a long peptide}

For a candidate minimal epitope $e_i$ and restricting HLA allele $h_i$, IRIS
produces a positive mechanistic score
\begin{equation}
s_i=F_i(e_i,h_i).
\end{equation}
The details of $F_i$ belong to Level~1. For Level~2, the important fact is
that a long peptide $\ell$ can contain a roster
\begin{equation}
\mathcal C_\ell=\{(e_i,h_i):i=1,\ldots,K_\ell\}
\end{equation}
of several candidate minimal-epitope--HLA pairs.

The observed PDAC and COVID response label belongs to $\ell$, not to an
identified member of $\mathcal C_\ell$. Level~2 must therefore map all
$K_\ell$ candidate scores to one long-peptide score $S_\ell$.

Because the Level~1 score is positive and multiplicative, it is convenient to
work on the log scale:
\begin{equation}
z_i=\log(s_i+\epsilon),
\end{equation}
where $\epsilon$ is the declared numerical floor. The transformation is
computational; neither $s_i$ nor $\sigmoid(z_i)$ is a response probability.
The Level~2 map developed below is therefore
\begin{equation}
S_\ell=S(z_1,\ldots,z_{K_\ell}).
\end{equation}

\subsection{The two biological limits}

Two plausible configurations motivate the aggregation problem.

\textbf{One-candidate sufficiency.}
One candidate may be strong enough to explain the long-peptide response.
Other roster members should not dilute it or automatically improve its score.

\textbf{Multi-candidate corroboration.}
Several candidates may jointly provide persuasive evidence even when no
single candidate is decisive. Their magnitudes should be allowed to
accumulate.

Maximum aggregation represents the first configuration:
\begin{equation}
S_{\max}(z)=\max_i z_i.
\end{equation}
It preserves the strongest candidate but discards every other score.

Log-sum-exp represents the second:
\begin{equation}
S_{\LSE}(z)=\log\sum_i e^{z_i}.
\end{equation}
It accumulates all candidates, but it can increase merely because a peptide
offers more candidate opportunities. The adaptive rule is designed to move
between these limits according to the score geometry of the individual
roster.

\subsection{What the original PDAC--COVID round robin showed}

The motivating comparison crossed five component constructions with twelve
fixed Level~2 aggregation rules, producing 60 complete systems. Its results
did not identify one aggregation behavior that was adequate in every disease
context.

\begin{itemize}
  \item In the threshold-zero COVID-only comparison, 11 systems were
  graph-maximal under PR and two under AUROC. The two common systems paired
  Mono-Q/full-PN with either full log-sum-exp or top-10 log-sum-exp. Within
  that component family, those two accumulation rules were supported over
  maximum, mean, median, log-mean-exp, the tested top-fraction means, and the
  tested top-$k$ means in both COVID datasets and both metrics.

  \item Original PDAC produced 14 PR-maximal and 13 AUROC-maximal systems,
  with ten common to both metrics. That set did not contain the COVID
  Mono-Q/full-PN log-sum-exp system and instead contained several
  log-mean-exp, top-fraction, top-$k$ mean, and top-$k$ log-sum-exp
  constructions.

  \item When PDAC, COVID spike, and COVID non-spike were required jointly, the
  formal frontier broadened to 44 PR-maximal and 42 AUROC-maximal systems,
  with 41 in their descriptive overlap. No system was separately maximal in
  all three contexts under either metric.
\end{itemize}

\paragraph{Assay geometry contextualizes the contrast.}
The aggregation rules did not face comparable candidate-opportunity sets.
COVID spike used 15-mers, and 1,564 of 1,607 COVID non-spike endpoints were
also 15-mers. PDAC used predominantly 27-mers: 182 of 213 endpoints had length
27, with a complete range of 15--49 amino acids. If candidate minimal epitopes
have lengths 9--12, a peptide of length $L\geq12$ contains
\begin{equation}
W(L)=\sum_{r=9}^{12}(L-r+1)=4L-38
\end{equation}
sequence windows before HLA expansion and scoreability filtering. Thus a
15-mer offers 22 windows and a 27-mer offers 70. The committed prediction
artifacts show that this difference persisted after filtering
(Table~\ref{tab:roster-geometry}). These are descriptive audits of the
original tournament inputs, not additional selection criteria.

\begin{table}[htbp]
\centering
\small
\caption{Long-peptide geometry and realized candidate rosters. $K$ counts
unique scoreable minimal-epitope--HLA tuples. Potential windows refer to the
predominant peptide length before HLA expansion and filtering; COVID labels
use the threshold-zero definition.}
\label{tab:roster-geometry}
\begin{tabularx}{0.98\textwidth}{@{}>{\raggedright\arraybackslash}Xrrrrr@{}}
\toprule
Cohort & Length & Windows & Median (mean) $K$ & $K=0$ & Mean $K$, neg./pos. \\
\midrule
COVID spike & 15 & 22 & $1\;(2.07)$ & 36.0\% & $1.30/2.48$ \\
COVID non-spike & mostly 15 & 22 & $3\;(3.46)$ & 32.7\% & $1.32/6.28$ \\
PDAC & mostly 27 & 70 & $10\;(10.84)$ & 0\% & $10.63/12.57$ \\
\bottomrule
\end{tabularx}
\end{table}

\paragraph{Sparse COVID rosters made accumulation selective.}
Fewer than one percent of the 3,372 COVID endpoints contained more than ten
candidates. Full and top-10 log-sum-exp therefore differed on only a small
fraction of endpoints, so their joint survival supports soft accumulation over
a generally sparse roster rather than a biologically distinguished cutoff at
ten. Conversely, because no COVID roster contained more than 15 candidates,
the tested 1\%, 2\%, and 5\% top-fraction means each selected one candidate and
were exactly score-vector-identical to maximum. Their separate graph vertices
do not represent distinct aggregation behavior. Candidate opportunity was
also aligned with the threshold-zero labels: positive endpoints carried more
candidates than negative endpoints, especially in non-spike, and
zero-candidate endpoints were concentrated among negatives. Among nonempty
rosters, the mean available LSE bonus $B$ was $0.517$ for non-spike positives
and $0.285$ for negatives; the corresponding spike values were $0.355$ and
$0.288$. Corroborative mass therefore tended to add more evidence to COVID
positives than to negatives.

\paragraph{Dense PDAC rosters made unrestricted accumulation less selective.}
Every PDAC endpoint contained at least one candidate, and 46\% contained more
than ten. The median available LSE bonus was $0.852$, compared with $0.325$ in
COVID non-spike and $0.230$ in COVID spike, but substantial bonus mass was
common to both PDAC classes: the mean was $0.910$ for positives and $0.791$
for negatives. Full log-sum-exp could therefore reward the broad candidate
opportunity supplied by a long peptide even when that breadth did not sharply
distinguish response. Maximum removed the explicit accumulation bonus but also
discarded potentially meaningful second- and third-ranked candidates; it does
not by itself eliminate opportunity effects on the extreme score. The PDAC
survival of top-fraction, top-$k$ mean, and top-$k$ log-sum-exp rules is
consistent with a compromise that retains restricted upper-tail corroboration
without accumulating the complete lower-scoring roster.

This interpretation separates Level~2 geometry from Level~1 component choice.
The Full-HLA and Mono-Q/full-PN constructions used the same candidate counts;
changing the source of $Q$ changed score magnitudes and ordering, not roster
size. Candidate density can therefore help explain why the cohorts favored
different aggregation behavior, but it cannot by itself explain the COVID
preference for Mono-Q/full-PN at a fixed aggregation rule. Nor do the observed
differences establish peptide length as their cause: the cohorts also differ
in disease, antigen source, immune history, assay design, HLA distribution,
and label construction.

The contrast nevertheless supplies the empirical motivation for an adaptive
rule. COVID supports corroborative accumulation where realized breadth is
sparse and response-associated, whereas PDAC prevents full log-sum-exp from
serving as a universal answer when candidate opportunity is dense and nearly
ubiquitous. Rather than choosing one fixed point between maximum and
log-sum-exp for every peptide, the proposed rule responds to the realized
leader, relative score mass, and candidate breadth of each roster. It also
lets admitted corroboration close its own gate as the aggregate becomes
stronger. This self-limitation addresses unrestricted accumulation without
treating cohort identity or peptide length as an input. It does not, by
itself, identify or remove every source of candidate-opportunity confounding.
The formal design requirements are stated after the rule is constructed.

\section{Constructing the adaptive rule}

The rule can be summarized in one sentence:

\begin{quote}
Start with the strongest candidate, form a score-weighted corroboration offer
from the roster, and admit only as much of that offer as the resulting
aggregate still needs.
\end{quote}

\subsection{Step 1: identify the leader and the available LSE bonus}

Let
\begin{equation}
m=\max_i z_i
\end{equation}
be the strongest candidate score. Express every candidate relative to that
leader:
\begin{equation}
r_i=e^{z_i-m},
\qquad
R=\sum_i r_i.
\end{equation}
At least one $r_i$ equals one, and all others lie in $(0,1]$.

The full bonus available between maximum and log-sum-exp is
\begin{equation}
B(z)=\log R.
\end{equation}
Indeed,
\begin{equation}
S_{\LSE}(z)=m+B(z).
\end{equation}
Thus $m$ is the maximum candidate score and $B$ is the available LSE bonus.

\subsection{Step 2: measure candidate breadth at a declared Hill order}

Normalize the relative strengths:
\begin{equation}
p_i=\frac{r_i}{R}.
\end{equation}
The $p_i$ sum to one and describe how the represented score mass is distributed
across candidates. They are normalized weights, not response probabilities.

For Hill order $q\in[0,\infty]$, define the effective candidate number
\begin{equation}
N_q(p)=
\begin{cases}
K, & q=0,\\[3pt]
\displaystyle
\exp\!\left(-\sum_i p_i\log p_i\right), & q=1,\\[7pt]
\displaystyle
\left(\sum_i p_i^q\right)^{1/(1-q)},
& q\notin\{0,1,\infty\},\\[7pt]
\displaystyle
\frac{1}{\max_i p_i}, & q=\infty.
\end{cases}
\label{eq:hill-number}
\end{equation}
Map it to the bounded breadth coefficient
\begin{equation}
D_q(z)=1-\frac{1}{N_q(p)}.
\label{eq:hill-breadth}
\end{equation}
For every order, $D_q=0$ for a singleton and $D_q=1-1/K$ when the
$K$ candidates have equal weights. The orders differ only for uneven rosters.
Hill numbers are nonincreasing in $q$: low orders recognize more of the roster
as breadth, whereas high orders require corroborative mass to lie closer to
the leader.

\begin{center}
\small
\begin{tabularx}{0.92\textwidth}{@{}p{0.12\textwidth}>{\raggedright\arraybackslash}X@{}}
\toprule
Order & Breadth interpretation \\
\midrule
$q=0$ & Richness-sensitive: every unique candidate contributes to $D_q$. \\
$q=1$ & Entropy-weighted: low-weight candidates contribute in proportion to
their information mass. \\
$q=2$ & Concentration-sensitive: $D_2=1-\sum_i p_i^2$. \\
$q=\infty$ & Leader-share-sensitive: $D_\infty=1-\max_i p_i$. \\
\bottomrule
\end{tabularx}
\end{center}

At $q=0$, the breadth coefficient depends on $K$, but the complete
corroborative quantity below does not depend on count alone: the available
bonus $B$ still weights candidates by their realized scores.

\subsection{Step 3: form the corroboration offer}

Define
\begin{equation}
C_q(z)=B(z)D_q(z).
\label{eq:corroboration-offer}
\end{equation}
The available bonus $B$ asks how much relative score mass the roster contains.
The coefficient $D_q$ asks how that mass should count as breadth at the
declared Hill order. Their product $C_q$ is the corroboration offered before
sufficiency gating.

\subsection{Step 4: define an aggregate-sufficiency gate}

For any provisional aggregate score $a$, define the decreasing gate
\begin{equation}
g_{c,\kappa}(a)
=
\sigmoid\!\left(\frac{c-a}{\kappa}\right)
=
\frac{1}{1+\exp[(a-c)/\kappa]},
\qquad \kappa>0.
\label{eq:aggregate-gate}
\end{equation}
The center $c$ is the aggregate score at which the gate equals one-half. The
width $\kappa$ controls how gradual the transition is: the gate falls from
$0.9$ to $0.1$ over an interval of length
$2\log(9)\kappa\approx4.39\kappa$.

A low aggregate leaves the gate open and permits corroboration. As admitted
corroboration raises the aggregate, the same gate progressively closes. A
strong leader starts the aggregate high and therefore makes the rule max-like
from the outset.

\subsection{Step 5: let the aggregate gate itself}

The self-gated adaptive Hill-$q$ aggregate is the unique solution of
\begin{equation}
\boxed{
S_\ell
=
\underbrace{m}_{\substack{\text{maximum}\\\text{candidate score}}}
+
\underbrace{C_q(z)}_{\substack{\text{corroboration}\\\text{offer}}}
\times
\underbrace{g_{c,\kappa}\!\left(S_\ell\right)}_
{\substack{\text{aggregate-sufficiency}\\\text{fraction}}}
}.
\label{eq:self-gated-aggregate}
\end{equation}
The equation is implicit by design. The gate is evaluated at the result after
corroboration has been admitted, rather than at the pre-aggregation maximum
candidate score $m$ before admission.
The rule therefore admits corroboration only while the resulting aggregate
still leaves room for it.

\begin{center}
\small
\begin{tabularx}{0.92\textwidth}{@{}>{\raggedright\arraybackslash}p{0.16\textwidth}>{\raggedright\arraybackslash}X@{}}
\toprule
Quantity & Interpretation \\
\midrule
$m$ & \textbf{Maximum candidate score.}  This is the score of the single
best candidate and therefore the evidence already available without invoking
any corroboration.  Starting from $m$ guarantees that additional candidates
can only add to, rather than dilute, the strongest signal; for a singleton,
the aggregate reduces exactly to $m$. \\
\addlinespace
$B$ & \textbf{Available LSE bonus.} Because
$B=S_{\LSE}-m=\log\sum_i e^{z_i-m}$, it is the entire gap between max pooling
and log-sum-exp pooling.  Candidates near the leader contribute substantially
to this offer, whereas candidates far below the leader contribute little.
Crucially, $B$ is only the amount \emph{available}; the rule does not accept
all of it automatically. \\
\addlinespace
$D_q$ & \textbf{Breadth fraction at Hill order $q$.} The order determines
how strongly low-weight candidates count as breadth. At $q=0$, every unique
candidate contributes to the coefficient; at larger $q$, consequential mass
must be increasingly concentrated near the leader. \\
\addlinespace
$C_q$ & \textbf{Corroboration offer.} The product $BD_q$ combines realized
relative score mass with the declared interpretation of breadth. It is the
largest bonus available before the aggregate-sufficiency gate is applied. \\
\addlinespace
$g_{c,\kappa}(S_\ell)$ & \textbf{Fraction permitted by aggregate sufficiency.}
This fraction is near one while the developing aggregate is low, equals
one-half at $S_\ell=c$, and approaches zero as $S_\ell$ becomes high. A large offer
therefore helps close its own gate rather than being accepted without limit. \\
\bottomrule
\end{tabularx}
\end{center}

\subsection{Why the implicit definition is safe}

Let $C=C_q(z)$ and define
\begin{equation}
H(a)=a-m-Cg_{c,\kappa}(a).
\end{equation}
If $C=0$, Equation~\eqref{eq:self-gated-aggregate} gives $S_\ell=m$. If $C>0$,
then
\begin{equation}
H(m)=-Cg_{c,\kappa}(m)<0,
\qquad
H(m+C)=C\{1-g_{c,\kappa}(m+C)\}>0.
\end{equation}
Moreover, because the gate is strictly decreasing,
\begin{equation}
H'(a)=1-Cg'_{c,\kappa}(a)>0.
\end{equation}
There is therefore exactly one solution in $[m,m+C]$. It follows immediately
that
\begin{equation}
m
\leq
S_\ell
\leq
m+C_q(z)
\leq
m+B(z)
=
S_{\LSE}(z).
\label{eq:self-gated-bounds}
\end{equation}
The score always lies between maximum and log-sum-exp. Increasing the
corroboration offer cannot lower the solution: implicit differentiation gives
\begin{equation}
\frac{\partial S_\ell}{\partial C}
=
\frac{g_{c,\kappa}(S_\ell)}{1-Cg'_{c,\kappa}(S_\ell)}
>0.
\end{equation}
This monotonicity is one reason to prefer the self-consistent gate over the
explicit shortcut $m+Cg_{c,\kappa}(m+C)$, whose accepted bonus can decrease
when the offer becomes larger.

The unique root can be computed deterministically by bisection on $[m,m+C]$.
The interval is known in advance, the root is scalar, and no naive fixed-point
iteration is required. Solver tolerance and iteration limits form part of the
scoring specification.

\subsection{Full design requirements and roster comparability}

The construction separates three questions. The available bonus $B$ asks how
much relative score mass is present. The Hill coefficient $D_q$ asks how the
distribution of that mass should count as breadth. The self-gate $g(S_\ell)$ asks
how much corroboration the resulting aggregate still needs. A complete
adaptive Level~2 design must also state how it treats differences in roster
construction without assuming that roster size is either automatically
supportive or automatically suspect.

It is useful to state the full design problem explicitly. An adaptive
long-peptide aggregate should solve, or at least expose, each of the following
problems.

\begin{enumerate}[label=\textbf{P\arabic*.},leftmargin=*]
  \item \textbf{Unresolved causation.}
  The response belongs to a long peptide, but the causal minimal
  epitope--HLA pair is not observed. The rule must aggregate the candidate
  roster without pretending that it has identified which candidate caused the
  response.

  \item \textbf{One-candidate sufficiency.}
  One strong candidate may be sufficient. Additional weak candidates should
  neither dilute that leader nor automatically create appreciable evidence.

  \item \textbf{Meaningful corroboration.}
  Several consequential candidates may jointly support a response even when
  no single candidate is decisive. A rule that always returns the maximum
  discards this information.

  \item \textbf{Dominance versus breadth.}
  A roster containing one leader and many negligible candidates should be
  distinguished from a roster in which several candidates carry comparable
  score mass. The available bonus $B$ records total relative mass, while $D_q$
  controls how its distribution is interpreted as breadth. At low $q$, more
  of the roster contributes; at high $q$, breadth must be concentrated nearer
  the leader. The rule should make this choice explicit rather than embedding
  one privileged order without examination.

  \item \textbf{Aggregate sufficiency.}
  The same corroboration offer need not have the same meaning at every
  absolute score level. A strong leader may need little help. A weak leader
  may admit corroboration, but the admitted amount should diminish as the
  resulting aggregate itself becomes strong. Gating only at the initial
  leader cannot provide this feedback.

  \item \textbf{Candidate-opportunity comparability.}
  Candidate rosters may differ in size because of peptide length, HLA coverage,
  or roster-construction choices. The role of roster richness must therefore
  be declared. At $q=0$, every unique tuple contributes to the breadth
  coefficient; at higher orders, low-weight tuples contribute less. In every
  case, $B$ still controls their realized score mass and the self-gate limits
  how far the offer moves the aggregate. The rule does not apply an additional
  penalty based only on $K$, and it does not claim to eliminate every
  opportunity effect caused by peptide length or enumeration depth.

  \item \textbf{Exact-duplicate control.}
  The aggregation operates on unique minimal-epitope--HLA tuples. If the same
  tuple appears more than once---for example, because a patient is homozygous
  for the HLA allele---it must enter the roster only once; otherwise repeated
  representation would manufacture breadth. By contrast, the same minimal
  epitope paired with different HLA alleles forms distinct tuples and remains
  as distinct candidates. The adaptive aggregation does not collapse or
  otherwise specially adjust those cross-HLA candidates.

  \item \textbf{Symmetry and singleton coherence.}
  Candidate order must not matter, tied candidates must be treated
  symmetrically, and a roster containing one candidate must return that
  candidate's Level~1 log score.

  \item \textbf{Auditability and label blindness.}
  Every adaptation term must be calculable from the candidate scores under the
  prespecified roster convention. Endpoint response labels must not choose the
  mixing behavior used for that endpoint.
\end{enumerate}

Problems P4 and P6 are especially easy to conflate. Hill breadth describes
how a declared order reads one realized weight distribution. Candidate-
opportunity comparability asks whether a difference in roster construction
represents additional biological evidence, harmless enumeration, or
duplicated evidence. The aggregation function exposes the first choice. It
cannot settle the second question without assumptions about how candidate
opportunities are generated.

The recommended rule therefore does \emph{not} contain a separate roster-size
penalty that lowers the aggregate as roster size increases. Such a penalty
would assign a lower score to the same meaningful score geometry merely
because more candidates were enumerated.
That ordering would be defensible only under an explicit stochastic model in
which additional candidate opportunities generate spurious favorable scores.
Without such a model, the penalty could remove genuine biological signal and
make the result depend on pipeline exhaustiveness.

Candidate opportunity enters through both the realized score mass and the
declared Hill order. If an additional candidate lies $d$ log-score units below
the leader, its relative weight is $e^{-d}$. Given the previous
relative-weight sum $R$, its increment to the available LSE bonus is
\begin{equation}
\log\!\left(R+e^{-d}\right)-\log R
=
\log\!\left(1+\frac{e^{-d}}{R}\right)
\leq
\frac{e^{-d}}{R}.
\end{equation}
Thus additional candidates have little effect through $B$ when their
\emph{combined} relative score mass is negligible. At $q=0$, however, adding a
new unique tuple also changes $D_0=1-1/K$; this richness sensitivity is an
explicit modeling choice rather than an accidental invariance claim. Exact-
duplicate control occurs before the score vector $z$ is formed: the roster is
a set, rather than a multiset, of minimal-epitope--HLA tuples. A homozygous
recurrence of the same tuple therefore contributes one candidate, whereas the
same minimal epitope paired with two different HLA alleles contributes two
distinct corroborative contributions. Overlapping but nonidentical minimal
epitopes likewise remain distinct. The formula does not assert that these
contributions are statistically independent.

The recommended aggregate consequently remains
$S_\ell$, with no separate roster-size penalty and no
cohort-identity input. The self-gate limits accumulation according to the
score the roster actually reaches; it should not be described as a complete
correction for roster opportunity.

\clearpage
\section{How the rule behaves}

\subsection{Limiting cases}

Equation~\eqref{eq:self-gated-bounds} guarantees that the result always lies
between maximum and log-sum-exp. The fixed point adds a second form of
adaptation: a large corroboration offer raises the aggregate, and that increase
progressively closes the same gate through which the offer is admitted.

\begin{center}
\small
\begin{tabularx}{0.96\textwidth}{@{}>{\raggedright\arraybackslash}p{0.28\textwidth}>{\raggedright\arraybackslash}p{0.22\textwidth}>{\raggedright\arraybackslash}X@{}}
\toprule
Candidate pattern & Offer and gate & Result \\
\midrule
One candidate & $C_q=0$ & exactly maximum \\
One dominant candidate & small $B$, hence small $C_q$ & approximately maximum \\
Weak sparse corroboration & modest offer; gate initially open & substantial partial accumulation \\
Weak dense corroboration & large offer closes its own gate & self-limited accumulation \\
Strong leader & aggregate starts high; gate closed & max-like behavior \\
\bottomrule
\end{tabularx}
\end{center}

Candidate order does not matter. A singleton returns its Level~1 log score
exactly. Tied candidates are treated symmetrically. For an empty roster, the
declared floor score is returned before Equation~\eqref{eq:self-gated-aggregate}
is evaluated.

When $\kappa$ approaches zero, the transition becomes nearly hard. Its
behavior is then approximately
\begin{equation}
S_\ell
\approx
\begin{cases}
m, & m\geq c,\\
m+C_q, & m+C_q\leq c,\\
c, & m<c<m+C_q.
\end{cases}
\label{eq:hard-gate-limit}
\end{equation}
Thus a strong leader receives essentially no bonus, a weak roster whose full
offer remains below $c$ can receive nearly all of it, and an offer that would
cross $c$ is admitted only until the developing aggregate reaches the gate's
transition region. For finite $\kappa$, these cases join smoothly rather than
forming a strict cap.

\clearpage
\subsection{Worked score geometries}

The following examples use $q=2$, $c=-1$, and $\kappa=0.5$ only to illustrate
the calculation. Equal-weight plateaus have the same Hill number at every
order, so changing $q$ does not alter those rows.

\begin{center}
\centering
\small
\textbf{Four illustrative candidate-score geometries.}

\begin{tabularx}{0.99\textwidth}{@{}>{\raggedright\arraybackslash}p{0.16\textwidth}>{\raggedright\arraybackslash}p{0.16\textwidth}rrrr>{\raggedright\arraybackslash}X@{}}
\toprule
Roster & $z$ & $D_2$ & $g(S_\ell)$ & Max & LSE & Self-gated result \\
\midrule
Isolated leader
& $(0,-5,-5)$ & $0.026$ & $0.119$ & $0$ & $0.013$
& $0.00004$: negligible corroboration \\
Weak plateau
& $(-2,-2,-2)$ & $0.667$ & $0.720$ & $-2$ & $-0.901$
& $-1.473$: broad evidence receives a substantial bonus \\
Strong plateau
& $(0,0,0)$ & $0.667$ & $0.104$ & $0$ & $1.099$
& $0.076$: the strong aggregate closes most accumulation \\
Dense weak plateau
& ten scores at $-2$ & $0.900$ & $0.491$ & $-2$ & $0.303$
& $-0.982$: the large offer raises the score toward $c$ and closes its gate \\
\bottomrule
\end{tabularx}
\end{center}

The two three-candidate plateaus have identical breadth and offers but
different absolute score levels. The dense weak plateau has a much larger
offer than the weak three-candidate plateau. Its aggregate nevertheless stops
near the gate center because admitted corroboration progressively reduces
$g(S_\ell)$. The calculation is endpoint-adaptive: no cohort label or peptide
length enters the equation.

\section{Parameter selection and evidential status}

The functional form is shared, but its parameters need not be identical for
Level~1 score constructions with different candidate-score geometries. Let
$u$ index the component model and $b\in\{\mathrm{PR},\mathrm{ROC}\}$ index the
metric branch. The recommended contract selects one triple
\begin{equation}
(q_{ub},c_{ub},\kappa_{ub})
\end{equation}
for each component--metric group. That triple is then fixed across original
PDAC, threshold-zero COVID spike, and COVID non-spike. The aggregation rule
therefore remains unaware of dataset identity: context adaptation occurs only
through the endpoint's $m$, $B$, $D_q$, and self-consistent aggregate.

Selection should evaluate one frozen joint Cartesian grid
\begin{equation}
\mathcal G=\mathcal Q\times\mathcal C\times\mathcal K
\end{equation}
for every component--metric group. Within each cohort, all triples are ranked
jointly by the branch metric. Cohorts receive equal weight. The primary choice
minimizes the equal-cohort mean fractional rank; declared tie-breaks minimize
the worst-cohort fractional rank and then mean metric regret before applying
deterministic numerical tie-breaks. PR selection uses ordinary empirical
average precision with threshold-block tie handling. ROC selection uses AUROC
with half credit for tied positive--negative pairs.

The entire joint $(q,c,\kappa)$ grid $\mathcal G$ must be frozen before the
surfaces are examined. Adding intermediate $q$, $c$, or $\kappa$ values after
viewing a rank surface changes a rank-based objective through joint-grid
density.
The output should retain the complete surfaces, the selected triple, the
number and composition of near-optimal triples, and whether $c$ or $\kappa$
lies on a numerical boundary. An endpoint Hill order $q=0$ or $q=\infty$ is a
scientifically interpretable member of the frozen candidate set, not an
instruction to extend the order grid.

Two diagnostics are especially important. First, report which Hill orders
occur among triples within a declared tolerance of the optimum, rather than
only the single lexicographic winner. Second, repeat the selection after
leaving out each cohort in turn. These analyses distinguish a stable
component--metric calibration from a triple whose apparent adequacy depends
on one context or a narrow $\kappa$ transition.

All endpoint adaptation is label blind once the triple is fixed. Response
labels participate only in the model-development selection of
$(q,c,\kappa)$; they never choose the gate or Hill order for an individual
endpoint. If the same PDAC and COVID cohorts later supply a directed
round-robin comparison of the selected scores, that evidence is conditional
post-selection evidence rather than independent confirmation. The parameter
selection, complete surfaces, solver contract, source checksums, and this
evidential status should be recorded in the selection manifest.

\section{Interpretive qualifications}

\textbf{Gate threshold.}
All candidates are represented on the shared Level~1 scale
$z_i=\log(F_i+\epsilon)$. Accordingly, $c$ is the aggregate log-$F$ score at
which the aggregate-sufficiency fraction equals one-half. It is no longer a
threshold applied only to the maximum candidate. The definition of $F$, the
logarithmic transformation, and $\epsilon$ must therefore remain fixed when
$c$ and $\kappa$ are specified.

\textbf{Hill order.}
The order $q$ controls how the normalized candidate-weight distribution is
translated into breadth. It is a calibration of the Level~2 rule, not a
probability of candidate independence. A model--metric-specific $q$ remains
fixed across cohorts after selection.

\textbf{Candidate-opportunity comparability.}
The rule contains no separate penalty based on roster size. Roster richness
can nevertheless enter through $D_q$, especially at $q=0$, while $B$ records
the candidates' combined relative score mass and the self-gate limits accepted
accumulation. This is a declared treatment of candidate opportunity, not proof
that peptide length or enumeration depth is harmless. Exact duplicate
minimal-epitope--HLA tuples are removed before aggregation, including repeated
occurrences caused by HLA homozygosity. Distinct overlapping epitopes and the
same epitope associated with different HLA alleles remain separate
corroborative contributions.

\textbf{Numerical solution.}
Equation~\eqref{eq:self-gated-aggregate} is solved as a scalar root on the
certified interval $[m,m+C_q]$. The bisection tolerance, maximum iteration
count, empty-roster floor, and behavior on any numerical failure must be fixed
and audited. A naive fixed-point iteration is not part of the definition.

\textbf{Scope of the revision.}
The self-gate limits an offered bonus according to the aggregate it produces.
It does not establish that the roster convention is biologically complete,
that distinct tuples are statistically independent, or that the selected
parameters will transfer to an unrepresented context. Those claims require
additional evidence.

\end{document}
